\begin{abstract}
Using the conformal and metric-measure companion results and the stated
Brownian/Liouville quantum gravity inputs, we prove spectral convergence
for critical spherical FK--Ising maps to Liouville Brownian motion on the
ordinary unit-area $\sqrt3$-quantum sphere. The discrete walk has total
attempt rate one, uses every map edge including loops and multiplicities,
and has the corner measure as its stationary law. Accelerating time by the
number of map edges gives joint convergence of the metric-measure space,
all ordered eigenvalues with multiplicities and padding, and the heat trace
locally uniformly at strictly positive times. The conductivity and clock
constant are one in these conventions.
\end{abstract}

\section{Introduction}\label{sec:introduction}

Random planar maps provide discrete random surfaces on which geometry and
statistical mechanics interact.  The random-cluster formulation of Fortuin
and Kasteleyn includes the Ising model at cluster weight two
\cite{FortuinKasteleyn1972RandomCluster}.  For maps decorated by the critical
Fortuin--Kasteleyn (FK) model, Sheffield's inventory accumulation bijection
gives an exact word encoding \cite{Sheffield2016Inventory}.  Its spanning-tree
specialization is the correspondence of Mullin studied by Bernardi, and
Sheffield identifies its fixed-map construction with Bernardi's subgraph
bijection \cite{Mullin1967TreeRooted,Bernardi2007TreeRooted,Bernardi2008Subgraphs,Sheffield2016Inventory}.
Sheffield also relates the infinite-volume inventory contours to correlated
Brownian motion \cite[Theorem~2.5]{Sheffield2016Inventory}.  The cone-time,
local, and finite-volume refinements of Gwynne, Mao, and Sun and of Gwynne
and Sun identify further contour and loop data
\cite{GwynneMaoSun2019ConeTimes,GwynneSun2017Local,GwynneSun2015Finite}.
The mating-of-trees theory of Duplantier, Miller, and Sheffield, and its
finite-sphere form due to Miller and Sheffield, identify the corresponding
Brownian sewing with a Liouville quantum gravity (LQG) surface explored by
space-filling
Schramm--Loewner evolution
\cite{DuplantierMillerSheffield2021Mating,MillerSheffield2019SphereMating}.
For the FK--Ising parameter \(q=2\), the resulting LQG parameter is
\(\gamma=\sqrt3\).  The Brownian paths in this encoding record the two
boundary-length contours; the vertex walk studied here is a separate
stochastic process on the map.

The continuum surface carries several related structures.  Its quantum area
comes from Gaussian multiplicative chaos (GMC) in the LQG construction of
Duplantier and Sheffield \cite{DuplantierSheffield2011KPZ}.  The construction of its intrinsic
metric proceeds through Liouville first-passage percolation: Ding, Dub\'edat,
Dunlap, and Falconet proved the required tightness, and Gwynne and Miller
proved existence, uniqueness up to a deterministic constant, and conformal
covariance of the LQG metric
\cite{DingDubedatDunlapFalconet2020Tightness,GwynneMiller2021Metric,GwynneMiller2021ConformalMetric}.
For discrete conformal coordinates, Gwynne, Miller, and Sheffield proved a
joint area and exploration limit for the Tutte embedding of disk mated-CRT maps,
including a Brownian walk limit modulo time parameterization
\cite{GwynneMillerSheffield2021Tutte}.  Holden and Sun proved joint embedded
metric and measure convergence for critical Boltzmann triangulations of polygons under
the Cardy embedding \cite{HoldenSun2023Cardy}.  These results illustrate the
importance of keeping metric, measure, and conformal coordinates coupled.

Liouville Brownian motion provides the diffusion associated with conformal
Brownian energy and quantum area.  Constructions were given by Berestycki and
by Garban, Rhodes, and Vargas
\cite{Berestycki2015Diffusion,GarbanRhodesVargas2016LBM}.  The Dirichlet-form
description and heat-kernel theory were developed further by Garban, Rhodes,
and Vargas, by Andres and Kajino, and by Maillard, Rhodes, Vargas, and
Zeitouni
\cite{GarbanRhodesVargas2014Dirichlet,AndresKajino2016Heat,MaillardRhodesVargasZeitouni2016Heat}.
In particular, the centered logarithmic Green operator gives a compact
spectral description in the torus setting of the last work.  Passing to the
ordinary unit-area quantum sphere requires a local comparison of its field
law with these Gaussian reference models.

For discrete walks, Berestycki and Gwynne proved convergence of walks on
mated-CRT maps to Liouville Brownian motion in whole-plane and disk settings,
with the limiting clock given a fixed annealed median exit normalization
\cite[Equation~(1.6) and Theorems~1.2 and~1.4]{BerestyckiGwynne2022MatedCRTLBM}.
Their deterministic discrete time normalization \(m_\varepsilon\) is defined by an
annealed median exit time and satisfies
\(m_\varepsilon\asymp\varepsilon^{-1}\); their Remark~1.3 leaves convergence
of \(\varepsilon m_\varepsilon\) unresolved.  Electrical estimates for
mated-CRT maps were established by Gwynne, Miller, and Sheffield
\cite{GwynneMillerSheffield2019Harmonic}, while Gwynne and Miller obtained
resistance, displacement, and spectral-dimension estimates for several
infinite planar-map ensembles \cite{GwynneMiller2021RandomWalk}.
To pass to a finite spectrum and heat trace, one must control the spatial
diffusion, its time normalization, and the spectral mass contributing at
positive times.

\subsection{The model and the theorem}

For \(n\ge1\), let \(\mathcal F_n\) be the finite set of pairs \((M,A)\),
where \(M\) is a connected planar multigraph with \(n\) edges embedded in the
oriented sphere and rooted at an oriented edge, and \(A\subseteq E(M)\).
Equivalence is by orientation-preserving homeomorphisms preserving the root
and \(A\).  Loops and multiple edges are allowed.  Let \(M^*\) be the dual
map and \(A^*=\{e^*:e\notin A\}\).  Write \(k(A)\) for the number of
components of the spanning subgraph \((V(M),A)\), including isolated
vertices, and set
\[
 \ell(M,A)=k(A)+k(A^*)-1=2k(A)+|A|-|V(M)|.
\]
The second identity is Euler's formula applied to the embedded spanning
subgraph.  The critical FK--Ising law is
\begin{equation}\label{eq:fk-law}
 \PP\bigl((M_n,A_n)=(M,A)\bigr)=\frac{2^{\ell(M,A)/2}}{Z_n},
 \qquad
 Z_n=\sum_{(M,A)\in\mathcal F_n}2^{\ell(M,A)/2}.
\end{equation}
Let \(V_n=V(M_n)\), let \(d_n\) be graph distance using every original map
edge with length one, and put
\begin{equation}\label{eq:corner-measure}
 \mu_n=\frac1{2n}\sum_{v\in V_n}\deg(v)\delta_v .
\end{equation}
A loop contributes two to the degree.  Thus \(\mu_n\) is normalized corner
counting pushed to vertices.

The walk waits an exponential time of mean one, chooses an incident half-edge
\(a\) uniformly, and moves to its opposite endpoint \(\operatorname{opp}(a)\).
A chosen loop leaves the position unchanged.  Its generator is
\begin{equation}\label{eq:walk-generator}
 (L_nf)(v)=\frac1{\deg(v)}
 \sum_{\substack{a\ \mathrm{half\text{-}edge}\\ a\ \mathrm{incident\ to}\ v}}
 \bigl(f(\operatorname{opp}(a))-f(v)\bigr).
\end{equation}
The operator \(-L_n\) is nonnegative and self-adjoint on \(L^2(\mu_n)\).
List its eigenvalues with multiplicity as
\[
 0=\lambda_{n,0}\le\lambda_{n,1}\le\cdots
 \le\lambda_{n,|V_n|-1},
\]
and put \(\lambda_{n,j}=+\infty\) for \(j\ge |V_n|\).  For \(t>0\), define
\begin{equation}\label{eq:discrete-heat-trace}
 H_n(t)=\Tr e^{ntL_n}
       =\sum_{j=0}^{|V_n|-1}e^{-nt\lambda_{n,j}}.
\end{equation}
This is the operator trace, with each mode counted once.

For the continuum object, let \((\Sph,h)\) be the ordinary unit-area
\(\gamma\)-Liouville quantum gravity sphere at \(\gamma=\sqrt3\):
the ordinary quantum-sphere measure of Duplantier, Miller, and Sheffield,
disintegrated at quantum area one, with its original marks forgotten
\cite[Definition~4.21(ii)]{DuplantierMillerSheffield2021Mating}.
Write \(\mu_h\) for its quantum area probability measure and \(D_h\) for
its intrinsic LQG metric in one fixed deterministic normalization.  We use
the conformal coordinate of the round sphere.  On smooth functions define
\begin{equation}\label{eq:continuum-energy}
 \mathcal E_h(f,g)=\frac12\int_{\Sph}
       \ip{\nabla f}{\nabla g}_{\mathrm{round}}\dd\mathrm{vol}_{\mathrm{round}}.
\end{equation}
The proof will establish that its closure on \(L^2(\mu_h)\), using the usual
quasi-continuous Sobolev representatives, is the Liouville Brownian motion
form and has compact resolvent.  Let \(A_h\) be its nonnegative operator and
\[
 0=\Lambda_0(h)\le\Lambda_1(h)\le\Lambda_2(h)\le\cdots
\]
its eigenvalues with multiplicity.  Set
\begin{equation}\label{eq:continuum-heat-trace}
 H_h(t)=\Tr e^{-tA_h}=\sum_{j\ge0}e^{-t\Lambda_j(h)},\qquad t>0.
\end{equation}
The factor \(1/2\) in \eqref{eq:continuum-energy} fixes the Brownian clock.
The two-dimensional Dirichlet integral is conformally invariant.  Pushing
\(\mu_h\) forward under a change of conformal coordinate therefore gives a
unitarily equivalent operator.

Here is the deterministic normalization of space.  For a finite real random
variable \(Y\), let \(\mathfrak m(Y)\) be the unique \(r\) for which
\(\EE[\arctan(Y-r)]=0\).  Put
\[
 \Delta_n=\max\{1,\diam(V_n,d_n)\},\qquad
 \Delta=\diam(\Sph,D_h),
\]
and define
\begin{equation}\label{eq:distance-normalization}
 a_n=\exp\{\mathfrak m(\log\Delta)-\mathfrak m(\log\Delta_n)\}.
\end{equation}
Both centers use unconditional laws.  In particular \(a_n\) is deterministic.
Lemma~\ref{lem:space-normalization} verifies that it agrees asymptotically
with the spatial normalization of the companion geometric limit.

\begin{theorem}[Joint Liouville spectral limit]\label{thm:main}
The form \eqref{eq:continuum-energy} is closable on \(L^2(\mu_h)\).
Its closure has compact resolvent, has the constants as its one-dimensional
kernel, and satisfies \(H_h(t)<\infty\) for every \(t>0\).
As \(n\to\infty\) through all positive integers,
\begin{equation}\label{eq:main-limit}
 \left(
   (V_n,a_nd_n,\mu_n),\
   (n\lambda_{n,j})_{j\ge0},\
   H_n
 \right)
 \ \Longrightarrow\
 \left(
   (\Sph,D_h,\mu_h),\
   (\Lambda_j(h))_{j\ge0},\
   H_h
 \right).
\end{equation}
The first coordinate has the Gromov--Hausdorff--Prokhorov topology on
isometry classes of compact metric probability spaces.  The second has the
product topology on \([0,\infty]^{\N_0}\).  The third has
\(C_{\mathrm{loc}}((0,\infty))\), the topology of uniform convergence on
each \([\delta,T]\) with \(0<\delta<T<\infty\).  All three limiting
coordinates are constructed from the same quantum sphere \(h\).
\end{theorem}

The clock is exactly \(n\) for the conventions above.  The reason is the
normalization of the electrical energy, together with a conductivity
identity proved below.  Formula~\eqref{eq:distance-normalization}, by
contrast, fixes only the spatial convention; the proof requires no power-law
formula for \(a_n\).

The proof uses the canonical conformal companion \cite{C} and the
metric-measure companion \cite{P}, the latter in its revision of
October 3, 2026.  Their encoding, protected local laws, quantitative address
and finite-view estimates, and geometric limits are the probabilistic inputs.
The electrical and spectral conclusions of
Theorem~\ref{thm:main} are obtained from the new arguments developed here.
In the companion titles, \emph{subcritical} refers to the LQG range
\(\gamma<2\); the FK model itself is at its self-dual critical point.

\subsection{Proof strategy}

There are two distinct questions.  The first is to identify the macroscopic
electrical behavior of the map.  The second is to control functions whose
variation is concentrated in a region carrying very little corner mass.
The first question determines the candidate limiting operator; the second
is needed to pass from that identification to the spectrum and heat trace.

We begin with the sum of squared voltage differences over the map edges.
Protected incidence records in the companions retain a finite graph together
with its ports and confinement.  We show that their local conditional-law
arguments also retain numerical electrical quantities of this graph, such
as effective resistance.  An adaptation of the aligned-annulus argument of
\cite[Section~8]{C} then rules out collapse of raw edge extremal length.
The key extra step compares the primal and dual networks in their common
flag surface.  A hypothetical collapse would produce cheap local crossings
in almost every position and direction for both networks.  Transverse
primal and dual crossings must use a paired edge, contradicting the
Cauchy--Schwarz inequality for their expected edge traffics.
Theorem~\ref{thm:network-bounds} turns this into cutoffs of bounded energy
and surrounding circuit flows of bounded squared traffic for fixed
geometric tests.

These fixed-test estimates identify the macroscopic energy.  We relax the
discrete energies under locally uniform convergence, using protected
finite-dimensional energy infima to compare inner and outer tests.  Cutoff
gluing makes the relaxation a local quadratic integral.  Circuit estimates
bound it above and below by the ordinary Sobolev energy.  Its matrix density
is a local germ, so the conditional independence and field-germ arguments
of \cite{C} make it a deterministic scalar.  The primal and dual scalars
agree by the exact FK duality.  In a square, their complementary
mixed-boundary conductances have product one; hence the scalar is one.
This uses the classical self-duality fixed-point idea of Dykhne
\cite{Dykhne1971}, after the existence, determinism, and equality of the two
scalar limits have been proved.  The result is
Theorem~\ref{thm:conductivity}.  A local minimization argument then
identifies the unit-current voltage between distinct macroscopic points
with the round Green voltage, yielding Theorem~\ref{thm:four-point}.

The remaining problem is quantitative compactness at smaller scales.
The address estimates of \cite{P} control the ranks of graphs formed by
exploration return relations.  We derive a version with one marked base
over a wider range of scales, and combine it with protected circuit flows
and a packing estimate for their reachable supports.  These circuits bound
oscillation on connected unions of their boundaries.  A bound on how long
the tour can remain inside one circuit disk lets us compare values associated
with parent and child blocks; telescoping these comparisons to singleton
blocks yields an approximation of all vertex values.

List the \(2n\) corner occurrences in inventory exploration order and
represent a vertex function \(f\) by its values on the corresponding \(2n\)
equal intervals of \([0,1]\).  This identifies \(L^2(\mu_n)\) isometrically with a subspace of
\(L^2([0,1])\).  Write \(\mathcal E_n^0(f)\) for the sum of squared differences
over all original edges, counted with their multiplicities.  After finitely
many charts are combined, Lemma~\ref{lem:energy-compactness} constructs
linear maps \(Q_{n,j}\) on these representatives, of rank at most \(C2^j\),
satisfying
\begin{equation}\label{eq:overview-approximation}
 \norm{f-Q_{n,j}f}_{L^2([0,1])}^2
 \le C2^{-pj}\mathcal E_n^0(f),\qquad p=\frac7{10}.
\end{equation}
On the joint representation extractions used below, almost surely there are
finite sample-dependent constants \(C,J,n_0\) for which both bounds hold
simultaneously for every index \(n\ge n_0\) in the extraction, every
\(j\ge J\), and every vertex function.  Theorem~\ref{thm:block-approximation}
proves the stronger local statement from which
\eqref{eq:overview-approximation} follows.  The estimate supplies both
strong compactness of bounded-energy functions and a polynomial lower bound
on high eigenvalues.

Section~\ref{sec:spectrum} combines the four-point voltage limit with this
approximation.  The former identifies double differences of the centered
inverse kernels away from coincident points.  Strong compactness prevents
an additional low mode from disappearing into a set of
vanishing mass, while \(p>1/2\) makes the inverse spectral tails square
summable.  An operator argument then upgrades the four-point identification
to Hilbert--Schmidt convergence of the centered inverses.  Ordered
eigenvalues and the heat traces follow.  This last passage is an abstract
compactness principle, related to the general theory of compact spectral
convergence of Kuwae and Shioya \cite{KuwaeShioya2003}: it uses finite-rank
energy approximations and four-point kernel limits, and does not depend on
the specific FK encoding.
The common canonical coordinate from \cite{C} keeps this spectral limit
coupled to the metric-measure limit of the same sphere.
